\documentclass[10pt,a4paper]{amsart}
\usepackage[margin=19mm]{geometry}
\usepackage{amsmath,amssymb,mathtools}
\usepackage[scr=rsfs,cal=euler]{mathalfa}
\usepackage{xfrac}
\usepackage[unicode,hidelinks,bookmarks=false]{hyperref}
\newcommand{\ie}{i.e.\ }
\newcommand{\cf}{cf.\ }
\newcommand{\resp}{respectively\ }
\newcommand{\Subsection}{\subsection}
\providecommand{\nobreakdash}{\nobreak\hbox{-}}
\setlength{\emergencystretch}{2em}
\multlinegap0pt
\usepackage{amssymb}
\usepackage{tikz}
\usepackage[shortlabels]{enumitem}
\newtheorem{lemma}{Lemma}
\newtheorem{claim}{Claim}
\usepackage{cleveref}
\crefname{lemma}{Lemma}{Lemmas}
\crefname{claim}{Claim}{Claims}
\newcommand{\N}{\mathbb N}
\newcommand{\dom}{\operatorname{dom}}
\newcommand{\im}{\operatorname{im}}
\newcommand{\makeset}[2]{\left\lbrace #1 \;\middle|\;
 \begin{tabular}{@{}l@{}}
   #2
  \end{tabular}
  \right\rbrace}
\title{The Zariski Topology on Homeomorphism groups}
\author{Luna Elliott}
\date{Source: arXiv:2601.15185v4 (CC BY 4.0)}
\begin{document}
\maketitle

\noindent\emph{Source and licence.} The Zariski Topology on Homeomorphism groups. Version arXiv:2601.15185v4 (CC BY 4.0), distributed by arXiv under the Creative Commons Attribution 4.0 International licence, http://creativecommons.org/licenses/by/4.0/, as recorded in the arXiv licence metadata for this exact version and shown on the abstract page https://arxiv.org/abs/2601.15185v4. Full licence notice: \texttt{licenses/arxiv-2601.15185-CC-BY-4.0.txt}.\par\smallskip

\noindent\textit{Editorial scope.} Counted unit: the article's lemma main (source0.tex lines 303--306). The article's complete original proof of it is delivered with it (source0.tex lines 307--413, one \texttt{proof} environment of 107 lines). No mathematics is written, repaired or re-derived: the statement and the proof are the article's own lines, in the article's own notation, and no retained line is substituted or re-flowed.  Retained uncounted as the source-local context the proof needs: lines 257--275; lines 260--265.  Nothing was omitted from any retained span, so the package carries no omission record.  No retained line contains a citation, so the wrapper carries no bibliography and no bibliography entry.  No retained line contains a figure environment, an image-inclusion command or drawing code, so no image is delivered and the package declares no asset dependency.  Editorial formatting only, in addition: an AMS \texttt{amsart} preamble that restates the article's own declarations and macros used by the retained lines, namely the environments \texttt{lemma}, \texttt{claim} on separate counters, together with the standard packages those lines need.  The retained environments print in this document as Lemma 1, Claim 1 in order of appearance, whereas the article prints the same items with the numbering of its own full document.  The complete native source of the article as distributed in the arXiv e-print for this exact version is bundled unchanged as \texttt{sources/193140/source0.tex}.
\begin{lemma}\label{lem:techZariski}
    Suppose that \(G\) is a group.
    Let \(\mathcal{B}_G\) denote the family of sets of the form
    \begin{equation}\label{bigeq}
       \makeset{x\in G}{\(1_G\neq x^{k_{0,0}}a_{0,1}\ldots a_{0,l_0-1}x^{k_{0,l_0-1}}a_{0,l_0}\)\\
   \(1_G\neq x^{k_{1,0}}a_{1,1}\ldots a_{1,l_1-1}x^{k_{1,l_1-1}}a_{1,l_1}\)\\
   \quad\quad\vdots\\
   \(1_G\neq x^{k_{m-1,0}}a_{m-1,1}\ldots a_{m-1,l_{m-1}-1}x^{k_{m-1,l_{m-1}-1}}a_{m-1,l_{m-1}}\)}
    \end{equation}
    where  
    \begin{itemize}
    \item \(m\in \N\), \(l_0,l_1,\ldots, l_{m-1}\in \N\backslash \{0\}\),
        \item  \(a_{i,j}\in G\) for \(i<m\) and \(0<j\leq l_i\),
        \item  \(a_{i,j}\neq 1_G\) for \(i<m\) and \(0<j<l_i\),
        \item  \(k_{i,j}\neq 0\) for \(i<m\) and \(0\leq j<l_i\).
    \end{itemize}
    The set \(\mathcal{B}_G\) is a basis for the Zariski topology on \(G\).
    Moreover, if \(\varnothing\notin \mathcal{B}_G\), then the Zariski topology is irreducible.
\end{lemma}

    \begin{equation}
       \makeset{x\in G}{\(1_G\neq x^{k_{0,0}}a_{0,1}\ldots a_{0,l_0-1}x^{k_{0,l_0-1}}a_{0,l_0}\)\\
   \(1_G\neq x^{k_{1,0}}a_{1,1}\ldots a_{1,l_1-1}x^{k_{1,l_1-1}}a_{1,l_1}\)\\
   \quad\quad\vdots\\
   \(1_G\neq x^{k_{m-1,0}}a_{m-1,1}\ldots a_{m-1,l_{m-1}-1}x^{k_{m-1,l_{m-1}-1}}a_{m-1,l_{m-1}}\)}
    \end{equation}

\begin{lemma}\label{main}
    Suppose that \(X\) is an infinite set and \(G\) is a highly transitive group of permutations of \(X\) such that every non-identity element has infinite support.
    In this case the Zariski topology on \(G\) is irreducible.
\end{lemma}

\begin{proof}
   Recall the basis \(\mathcal{B}_G\) from \cref{lem:techZariski}.
Let \(m\), \(l_i\), \(a_{i,j}\) and \(k_{i,j}\) and \(S\in \mathcal{B}_G\) be such that \(S\) is the set described in  \cref{bigeq}.
    By \cref{lem:techZariski}, we need only show that \(S\neq \varnothing\).

    We will use many partial permutations of \(X\) in building up to an element of \(S\), that is to say bijections between subsets of \(X\).
     We compose and invert these partial permutations as binary relations and denote the composition by concatenation.
     Let \(p_0,p_1,\ldots, p_{m-1}\in X\) be fixed and distinct.
    For a partial permutation \(\sigma\) of \(X\), we define \(d_\sigma:\{0,1,\ldots,m-1\}\to \N\) by
     defining \((i)d_{\sigma}\) to be the largest value of \(j\leq l_i\) such that \(p_i\in \operatorname{dom}(\sigma^{k_{i,0}}a_{i,1}\ldots \sigma^{k_{i,j-1}}a_{i, j})\) (when \(j=0\), we interpret \(\sigma^{k_{i,0}}a_{i,1}\ldots \sigma^{k_{i,j-1}}a_{i, j}\) as the identity with domain \(X\)). 
     For each partial permutation \(\sigma\) and \(i<m\), let 
     \[f_{i,\sigma}:=\sigma^{k_{i,0}}a_{i,1}\ldots \sigma^{k_{i,(i)d_\sigma-1}}a_{i, (i)d_\sigma}.\]
     In particular when \((i)d_{\sigma}=0\), the map \(f_{i,\sigma}\) is the identity function on \(X\).
     We also define
     \[B_\sigma :=\dom(\sigma)\cup \im(\sigma)\cup \{p_0,p_1,\ldots,p_{m-1}\}\cup \{(p_0)f_{0,\sigma},(p_1)f_{1,\sigma},\ldots, (p_{m-1})f_{m-1,\sigma}\}.\]
     Note that \(B_\sigma\) is always finite when \(\sigma\) is finite.
     We inductively build a finite partial permutation \(\sigma\) of \(X\) such that at each stage of the induction, the following conditions hold:
     \begin{enumerate}
      \item[($\sigma$1)] if \(i\neq i'<m\), then \((p_i)f_{i,\sigma}\neq (p_{i'})f_{i',\sigma}\),
         \item[($\sigma$2)] if \(i<m\) and \((i)d_\sigma \neq l_i\), then  \((p_i)f_{i,\sigma}\notin \dom(\sigma)\cup \im(\sigma)\),
         \item[($\sigma$3)] if \(i<m\) and \((i)d_\sigma\neq 0\), then  \((p_i)f_{i,\sigma}\neq p_i\).
     \end{enumerate}
     Moreover, at each stage of the induction we will increase one of the image points of \(d_\sigma\) (which are bounded by the \(l_i\)) so that the induction must eventually terminate. 
     At the end of the induction, we will have \((i)d_\sigma=l_i\) for all \(i<m\). 
     Once we have done this, we can extend \(\sigma\) to an element of \(x\in G\) due to \(G\) being highly transitive. Moreover, from the last condition of the induction, we will have \(p_i\neq (p_i)f_{i,\sigma}=(p_i)f_{i,x}\) for all \(i<m\). In particular we will have \(f_{i,x}\neq 1_G\) for all \(i<m\), so \(x\in S\) as required.

     It remains only to inductively build \(\sigma\).
     Initialize \(\sigma\) as the empty function. In particular \(d_\sigma\) is the constant zero function, each \(f_{i,\sigma}\) is the identity function and \(B_{\sigma}=\{p_0,p_1,\ldots, p_{m-1}\}\).
     
     For the inductive step, choose \(i<m\) such that \((i)d_{\sigma}<l_i\). We keep this \(i\) fixed for the remainder of the proof. The next claim helps us extend \(\sigma\).
     \begin{claim}
         There is a finite set \(C_\sigma\subseteq X\) and a finite partial bijection \(\gamma\) of \(X\) satisfying the following conditions:
         \begin{itemize}
             \item \(B_\sigma\subseteq C_\sigma\), 
             \item \((\dom(\gamma)\cup \im(\gamma))\cap B_{\sigma}=\{(p_i)f_{i,\sigma}\}\),
             \item \((p_i)f_{i,\sigma}\in \dom(\gamma^{k_{i,(i)d_{\sigma}}})\), 
             \item \((\dom(\gamma)\cup \im(\gamma))\backslash C_\sigma=\{(p_i)f_{i,\sigma}\gamma^{k_{i,(i)d_{\sigma}}}\}\), 
             \item \(\dom(\gamma)\cap \dom(\sigma)=\varnothing\),
             \item  \(\im(\gamma)\cap \im(\sigma)=\varnothing\),
             \item \((p_i)f_{i,\sigma}\gamma^{k_{i,(i)d_{\sigma}}}a_{i, (i)d_{\sigma}+1}\notin C_\sigma\)
             \item if \((i)d_{\sigma}+1\neq l_i\), then \((p_i)f_{i,\sigma}\gamma^{k_{i,(i)d_{\sigma}}}a_{i, (i)d_{\sigma}+1}\neq  (p_i)f_{i,\sigma}\gamma^{k_{i,(i)d_{\sigma}}}\).
         \end{itemize} 
     \end{claim}
     \begin{proof}[Proof of Claim]
     We define \(C_\sigma\) to be the set \(\{c_0,c_1,\ldots, c_{|k_{i,(i)d_{\sigma}}|-1}\}\cup B_{\sigma}\) where \(c_0=(p_i)f_{i,\sigma}\) and the points \(c_1,c_2,\ldots,c_{|k_{i,(i)d_{\sigma}}|-1}\in X\backslash B_\sigma\) are arbitrary but distinct.
     We also define another point \(c_{|k_{i,(i)d_{\sigma}}|}\) depending on the following case analysis:
     \begin{enumerate}
         \item[(i)] If \((i)d_\sigma\neq l_i-1\): 
         Choose \(c_{|k_{i,(i)d_{\sigma}}|}\in X\backslash C_\sigma\) such that \((c_{|k_{i,(i)d_{\sigma}}|})a_{i,(i)d_{\sigma}+1}\in X\backslash C_\sigma\) and  \((c_{|k_{i,(i)d_{\sigma}}|})a_{i,(i)d_{\sigma}+1}\neq c_{|k_{i,(i)d_{\sigma}}|}\). This is possible as \(a_{i,(i)d_{\sigma}+1}\) has infinite support and \(C_\sigma\) is finite.
       
       \item[(ii)] If \((i)d_\sigma= l_i-1\): 
         Choose \(c_{|k_{i,(i)d_{\sigma}}|}\in X\backslash C_\sigma\) such that \((c_{|k_{i,(i)d_{\sigma}}|})a_{i,l_i}\in X\backslash C_\sigma\). In this case it is possible that \(a_{i,(i)d_{\sigma}+1}\) is the identity function so we cannot make the same assumption as in case (i).
     \end{enumerate}
     If \(k_{i,(i)d_{\sigma}}>0\), we define \(\gamma: \{c_0,\ldots, c_{|k_{i,(i)d_{\sigma}}|-1}\}\to X\) by \((c_j)\gamma=c_{j+1}\). If \(k_{i,(i)d_{\sigma}}<0\), we define  \(\gamma: \{c_1,\ldots, c_{|k_{i,(i)d_{\sigma}}|}\}\to X\) by \((c_j)\gamma=c_{j-1}\).   

     We need to check that the choice of \(C_\sigma\) and \(\gamma\) satisfy the required conditions. \(B_\sigma \subseteq C_\sigma\) by definition. Note that \(\dom(\gamma)\cup \im(\gamma)=\{c_0, c_1,\ldots, c_{|k_{i,(i)d_{\sigma}}|}\}\) in both cases for \(\gamma\). The only element of this set not explicitly chosen to not belong to \(B_\sigma\) is \(c_0\). We also have \(c_0=(p_i)f_{i,\sigma}\in B_\sigma\). Thus \((\dom(\gamma)\cup \im(\gamma))\cap B_{\sigma}=\{(p_i)f_{i,\sigma}\}\). Also, by ($\sigma$2), it follows that \(\dom(\gamma)\cap \dom(\sigma)=\varnothing\) and \(\im(\gamma)\cap \im(\sigma)=\varnothing\).
     By the choice of \(\gamma\), \(\gamma^{k_{i,(i)d_{\sigma}}}\) has domain \(\{c_0\}=\{(p_i)f_{i,\sigma}\}\) and \((c_0)\gamma^{k_{i,(i)d_{\sigma}}}=c_{|k_{i,(i)d_{\sigma}}|}\).
     The only element of \(\dom(\gamma)\cup \im(\gamma)=\{c_0, c_1,\ldots, c_{|k_{i,(i)d_{\sigma}}|}\}\) which does not belong to \(C_\sigma\) is \(c_{|k_{i,(i)d_{\sigma}}|}=(c_0)\gamma^{k_{i,(i)d_{\sigma}}}=(p_i)f_{i,\sigma}\gamma^{k_{i,(i)d_{\sigma}}}\).
     In both case (i) and case (ii), the element \((c_{|k_{i,(i)d_{\sigma}}|})a_{i, (i)d_{\sigma}+1}=(p_i)f_{i,\sigma}\gamma^{k_{i,(i)d_{\sigma}}}a_{i, (i)d_{\sigma}+1}\) was chosen not to belong to \(C_\sigma\). Finally, if \((i)d_{\sigma}+1\neq l_i\), then from case (i) we also have \((p_i)f_{i,\sigma}\gamma^{k_{i,(i)d_{\sigma}}}=c_{|k_{i,(i)d_{\sigma}}|}\neq (c_{|k_{i,(i)d_{\sigma}}|})a_{i, (i)d_{\sigma}+1}=(p_i)f_{i,\sigma}\gamma^{k_{i,(i)d_{\sigma}}}a_{i, (i)d_{\sigma}+1}\) as required.
     \end{proof}
     
     Let \(C_\sigma\) and \(\gamma\) be as in the claim.
     By the claim, the domains and images of \(\sigma\) and \(\gamma\) are disjoint. Thus we can define a new finite partial permutation \(\beta\) of \(X\) with domain \(\dom(\sigma)\cup \dom(\gamma)\) by
     \[(p)\beta = \begin{cases}
         (p)\sigma & \text{if } p\in\dom(\sigma)\\
         (p)\gamma & \text{if } p\in\dom(\gamma).
     \end{cases}\]
     We need only show that \(\beta\) can serve as \(\sigma\) for the next step of the induction. That is to say, we need only check that the following conditions hold:
\begin{enumerate}
\item[($\beta$0)]  \((i)d_{\beta}>(i)d_{\sigma}\),
      \item[($\beta$1)] if \(j\neq j'<m\), then \((p_j)f_{j,\beta}\neq (p_{j'})f_{j',\beta}\),
         \item[($\beta$2)] if \(j<m\) and \((j)d_\beta \neq l_j\), then  \((p_j)f_{j,\beta}\notin \dom(\beta)\cup \im(\beta)\),
         \item[($\beta$3)] if \(j<m\) and \((j)d_\beta\neq 0\), then  \((p_j)f_{j,\beta}\neq p_j\).
     \end{enumerate}
     
     By construction, the point
     \((p_i)f_{i,\sigma}\beta^{k_{i,(i)d_{\sigma}}}=(p_i)f_{i,\sigma}\gamma^{k_{i,(i)d_{\sigma}}}\) is the unique element of the set \((\dom(\beta)\cup \im(\beta)) \backslash C_\sigma\). In particular \((i)d_\beta\geq (i)d_\sigma+1\). If \((i)d_\sigma=l_i-1\), then we must have \((i)d_\beta=(i)d_\sigma+1\). On the other hand, if \((i)d_\sigma\neq l_i-1\), then the claim shows that \[(p_i)f_{i,\sigma}\beta^{k_{i,(i)d_{\sigma}}}a_{i,(i)d_\sigma+1}=(p_i)f_{i,\sigma}\gamma^{k_{i,(i)d_{\sigma}}}a_{i,(i)d_\sigma+1}\] does not belong to \(C_\sigma\) and is not equal to \((p_i)f_{i,\sigma}\beta^{k_{i,(i)d_{\sigma}}}\). As \((p_i)f_{i,\sigma}\beta^{k_{i,(i)d_{\sigma}}}\) is the only element of \((\dom(\beta)\cup \im(\beta)) \backslash C_\sigma\), the point \((p_i)f_{i,\sigma}\beta^{k_{i,(i)d_{\sigma}}}a_{i,(i)d_\sigma+1}\) belongs to neither the domain nor the image of \(\beta\). 
     So 
     \begin{equation}\tag{$*$}\label{eq:d+1}
         (i)d_\beta=(i)d_\sigma +1
        \end{equation}in all cases. We have now established ($\beta$0). It remains only to check conditions ($\beta$1), ($\beta$2) and ($\beta$3).
     We have also just established that if  \((i)d_\beta \neq l_i\), then 
     \begin{equation}\tag{$**$}\label{eq:newendnowhere}
         (p_i)f_{i,\beta}=(p_i)f_{i,\sigma}\beta^{k_{i,(i)d_{\sigma}}}a_{i,(i)d_\sigma+1}\notin \dom(\beta)\cup \im(\beta)\cup C_{\sigma}.
     \end{equation}

     By the claim, the only element of \((\dom(\gamma)\cup \im(\gamma))\cap B_\sigma \) is the point \((p_i)f_{i,\sigma}\).
     By ($\sigma$1) and ($\sigma$2), we have \((p_{i'})f_{i',\sigma}\in B_\sigma\backslash (\dom(\sigma)\cup \im(\sigma))\) and \((p_i)f_{i,\sigma}\neq (p_{i'})f_{i',\sigma}\) for \(i'\neq i\).
     It follows that \((p_{i'})f_{i',\sigma}\not\in \dom(\beta)\cup \im(\beta)\) when \(i'\neq i\). Thus when \(i'\neq i\), we have 
     \begin{equation}\tag{$***$}\label{eq:otheris}
     (i')d_\sigma =(i')d_\beta\text{ and }(p_{i'})f_{i',\beta}=(p_{i'})f_{i',\sigma}\in B_{\sigma}.
     \end{equation}
     From \cref{eq:d+1} and \cref{eq:otheris}, we have now entirely determined the map \(d_{\beta}\) and the partial map \(f_{i,\beta}\).
     From the claim we have \((\dom(\gamma)\cup \im(\gamma))\cap B_{\sigma}=\{(p_i)f_{i,\sigma}\}\). This together with \cref{eq:otheris} and ($\sigma$2) show that \((p_{i'})f_{i',\beta}=(p_{i'})f_{i',\sigma}\notin \dom(\beta)\cup \im(\beta)\) for \(i\neq i'\) with \((i')d_\sigma\neq l_{i'}\).
     By \cref{eq:newendnowhere}, we also have (\(p_{i})f_{i,\beta}\notin \dom(\beta)\cup \im(\beta)\) when \((i')d_\sigma\neq l_{i'}\). We have now shown that condition ($\beta$2) holds. 
     Only conditions ($\beta$1) and ($\beta$3) remain.
     
      By \cref{eq:d+1} and the claim, we have \((p_i)f_{i,\beta}\notin C_\sigma\supseteq B_\sigma\). Thus, by \cref{eq:otheris}, \((p_i)f_{i,\beta}\) is distinct from each of \((p_{i'})f_{i',\beta}=(p_{i'})f_{i',\sigma}\) with \(i'\neq i\). The points \((p_{i'})f_{i',\beta}=(p_{i'})f_{i',\sigma}\) for various \(i'\neq i\) are also distinct by \cref{eq:otheris} and ($\sigma$1). Together these establish that condition ($\beta$1) holds. Only condition ($\beta$3) remains.

     Suppose now that \(i'<m\) and \((i')d_\beta \neq 0\).
     There are two cases. If \(i'\neq i\), then by \cref{eq:otheris} we have \((i')d_\sigma=(i')d_\beta \neq 0\) and \((p_{i'})f_{i',\beta}=(p_{i'})f_{i',\sigma}\).
     By ($\sigma$3), it follows that  \((p_{i'})f_{i',\beta}\neq p_{i'}\) as required.
     The remaining case is when \(i'=i\). By \cref{eq:d+1} and the claim,
     the point \((p_i)f_{i,\beta}\) does not belong to \(C_\sigma\).
     In particular, \((p_i)f_{i,\beta}\neq p_i\in B_\sigma\subseteq C_\sigma\).
     
\end{proof}

\end{document}
